% ============================================================================
% 摘要与关键词
% ============================================================================
\begin{center}
  \begin{minipage}{0.92\textwidth}
    \small
    \noindent{\sffamily\bfseries 摘\quad 要：}
    针对自动化立体仓库中多取货任务并发场景下的多穿梭无人车（AGV）智能调度问题，本文围绕“动态任务指派”与“无碰撞时空路径规划”的强耦合难题，设计并实现了一套集建模仿真、启发式调度优化与独立物理验解于一体的调度系统。系统采用“任务指派层—时空路径规划层—独立物理验解层”三层协同架构，在严格遵守货架层序解锁、窄巷道顶点与换边避碰、装卸占位及流关闭返库等物理约束的前提下，围绕通用多工况调度引擎（\code{src/agv/}）与真实工业订单流（WEPAStacks）两条互补主线开展系统研究与算法演进。在核心路径层，采用基于安全时间区间压缩的 SIPP 算法并结合增量预约表索引，显著降低了长时域避碰搜索开销；在决策层，面向一般分层与混合货架场景发展了拥堵关联大邻域修复（LNS）及高积压影子试派机制（BCRD），面向大规模真实订单流收敛出“事件门控滚动派单＋软分区联合遗传算法（GA）＋预约表 SIPP 物理预演”的综合方案（\code{gate}）。在涵盖 561 次独立验解运行的跨窗口多随机种子实验中，所有输出轨迹均 100\% 通过独立物理校验；在真实订单流既有验证窗口上，最终方案 \code{gate} 相对先到先服务（\code{fcfs}）、最早到达优先（\code{eta}）与滚动匈牙利指派（\code{hungarian}）分别将综合评价指标降低 5.45\%、2.43\% 与 3.87\%，并通过事件门控将开发窗 GA 触发次数削减 63\%。本文系统阐述了问题背景、方法框架与思路演进脉络，并为后续形式化建模、算法细节剖析与完整实验论证奠定结构基础。

    \vspace{5pt}
    \noindent{\sffamily\bfseries 关键词：}
    穿梭无人车（AGV）；多智能体路径规划（MAPF）；安全区间路径规划（SIPP）；联合遗传算法；事件门控滚动调度；立体仓库仿真
  \end{minipage}
\end{center}
\vspace{6pt}
\noindent\rule{\textwidth}{0.4pt}
\vspace{4pt}
